\subsection{Banach manifolds and Lie groups}\label{subse:banmf}

We need some background on (real) infinite-dimensional $C^p$ or smooth manifolds, modeled on Banach spaces, and the resulting notion of (Banach--)Lie group. There is no shortage of good sources in the literature, but they occasionally differ on the precise setup (e.g., in considering smooth vs.\ {\it analytic} manifolds).

\cite{bourb-vars-17,bourb-vars-8-15} give a very convenient uniform treatment of $C^p$ manifolds, where $p$ can be
\begin{itemize}\itemsep=0pt
\item a non-negative integer, whereby the $C^p$ condition (on a map, say) is defined recursively, as being differentiable with derivatives of class $C^{p-1}$ (and with $C^0$ signifying continuity);

\item the symbol `$\infty$', meaning of class $C^p$ for all non-negative integers $p$ (the term `smooth' is also customary for such manifolds/maps);

\item the symbol `$\omega$', meaning `analytic' (real, or complex, or over an ultrametric field).
\end{itemize}
As noted in \cite[Avertissement]{bourb-vars-17}, however, the work constitutes a summary of results and does not contain proofs. Lie {\it groups} are studied in a much more detailed fashion in \cite[Chapter III]{bourb-lie-13}, with the caveat that the text refers back to \cite{bourb-vars-17,bourb-vars-8-15} frequently.

Banach manifolds also receive a thorough treatment in \cite{lang-fund}, which focuses on the {\it smooth} (rather than analytic) setting. In particular, the Lie groups of \cite[Section~VI.5]{lang-fund} are required to be only ($C^p$ or) smooth; the same goes for those of \cite[Definition IV.1]{neeb-inf}. One the other hand, the Banach manifolds \cite[Section~3]{upm-ban} and Lie groups of \cite[Section~6]{upm-ban} are analytic.

Below, we work mostly with analytic (rather than smooth) manifolds and especially Lie groups. The ground field can be either $\bR$ or $\bC$, and the discussion will go through unaffected regardless of the choice. Over $\bC$ the analytic (i.e., holomorphic) setting is the only reasonable choice, while for real Banach {\it Lie groups} are automatically analytic \cite[Corollary IV.1.10]{neeb-lc}.

Natural examples of Lie groups include
\begin{itemize}\itemsep=0pt
\item Banach spaces with their additive structure \cite[Section~III.1.1, Examples following Proposition 3]{bourb-lie-13}, acting also as their own Lie algebras;
\item The multiplicative group $A^{\times}$ of invertible elements in a unital Banach algebra (\cite[Proposition IV.9]{neeb-inf} or \cite[Example 6.9]{upm-ban}, where the group is denoted by $G(A)$). The Lie algebra is $A$;
\item The unitary group $U(A)$ of a $C^*$-algebra $A$, with the space
 \begin{equation*}
 \{a\in A\mid a+a^*=0\}\subset A
 \end{equation*}
 of {\it skew-adjoint} elements as its Lie algebra \cite[Corollary~15.22]{upm-ban},
\end{itemize}
and so on.

\begin{Convention}\label{cv:qlie}
 Submanifolds $N\subseteq M$ and Lie subgroups (which will typically be closed), unless specified otherwise, are as in \cite[Section~5.8.3]{bourb-vars-17} (and \cite[Section~III.1.3, Definition~3]{bourb-lie-13}, \cite[Section~II.2]{lang-fund}, etc.): the requirement is that the resulting embeddings
 \begin{equation*}
 T_pN\le T_pM,\qquad p\in N
 \end{equation*}
 {\it split} as Banach-space morphisms, i.e., that there be closed subspaces
 \begin{equation*}
 X\le T_pM,\qquad T_pM\cong T_pN\oplus X.
 \end{equation*}

 By contrast, absent the splitting, the respective terms employed in \cite[Section~III.1.3, Remark]{bourb-lie-13} are {\it quasi-submanifold} and {\it Lie quasi-subgroup}. We use these too, on occasion.
\end{Convention}

\begin{Remark}
 These conventions are {\it not} universal: \cite[Definition IV.6]{neeb-inf}, for instance, requires closure of Lie subgroups, but not splitting. So,
 \begin{itemize}\itemsep=0pt
 \item the Lie quasi-subgroups of \cite[Section~III.1.3, Remark]{bourb-lie-13} are the (plain) Lie subgroups of \cite[Definition IV.6]{neeb-inf};
 \item while the Lie subgroups of \cite[Section~III.1.3, Definition 3]{bourb-lie-13} are termed {\it complemented} Lie subgroups in \cite[Definition IV.6]{neeb-inf}.
 \end{itemize}
\end{Remark}

We always assume Lie algebras $\fu:=\operatorname{Lie}(U)$ of Banach Lie groups to be equipped with norms making them into {\it Banach Lie algebras} \cite[discussion following Definition IV.1]{neeb-inf}: $(\fu,\|\cdot\|)$ is a~Banach space, and
\begin{equation*}
 \|[x,y]\|\le C\|x\| \|y\|,\qquad \forall x,y\in \fu \quad\text{for some}\quad C>0.
\end{equation*}

%%%%%%%%%%%%%%%%%%%%%%%%%%%%%%%%%%%%%%%%%%%%%%%%%%%%%%%%%%%%%%%%%%%%%%%%%%%%%
%%%%%%%%%%%%%%%%%%%%%%%%%%%%%%%%%%%%%%%%%%%%%%%%%%%%%%%%%%%%%%%%%%%%%%%%%%%%%
\section{Cocycle manifolds and discrete cohomology}\label{subse:cocyman}

In the discussion below, we will frequently have to work with spaces
\begin{equation*}
 C(X,M)
 :=
 \text{continuous maps }X\to M
\end{equation*}
for compact (always Hausdorff) spaces $X$ and Banach manifolds $M$. They will always be equipped with the {\it uniform topology} (or {\it compact-open topology} \cite[Definition preceding Theorem 46.8]{munk}, since $X$ is compact). For Banach {\it spaces} $E$, this topology on $C(X,E)$ is the one induced by the supremum (complete) norm
\begin{equation*}
 \|f\|:=\sup_{x\in X}\|f(x)\| = \max_{x\in X}\|f(x)\|,\qquad \forall f\in C(X,E).
\end{equation*}

As explained in \cite[Section~3.2]{ps-loop}, for a Banach Lie group $U$ with Lie algebra $\fu$ the space $C(X,U)$ is again a Banach Lie group, with $C(X,\fu)$ as its Lie algebra.

%%%%%%%%%%%%%%%%%%%%%%%%%%%%%%%%%%%%%%%%%%%%%%%%%%%%%%%%%%%%%%%%%%%%%%%%%%%%%
\subsection{Group cohomology (possibly non-abelian)}\label{subse:coh}

The setup will be that of compact (occasionally finite) groups $G$ acting on Banach Lie groups~$U$ so as to preserve the relevant structure (i.e., by analytic group automorphisms). Actions are always assumed continuous in the strongest reasonable sense:
\begin{equation*}
 G\times U\to U
\end{equation*}
is continuous. One frequently considers weaker notions of continuity (e.g., \cite[Section~X.1]{tak2}), but those will not play a role below.

For finite $G$, the reader can find quick treatments of the pertinent material on the resulting
\begin{itemize}\itemsep=0pt
\item {\it cohomology groups} $H^n(G,U)$, $n\in \bZ_{\ge 0}$ for abelian $U$;
\item and $H^n(G,U)$, $n=0,1$ for arbitrary $U$
\end{itemize}
in, say, \cite[Chapter VII and its Appendix]{ser-lf} or \cite[Sections~I.2 and I.5]{ser-gal}.

Briefly, the non-abelian machinery is as follows. Given an action
\begin{equation}\label{eq:gactu}
 G\times U\ni (g,u)\xmapsto{\quad} {}^gu\in U
\end{equation}
of a group $G$ (typically finite, for us) on a Lie group $U$ (possibly non-commutative and frequently infinite-dimensional), the {\it $U$-valued cocycles} attached to the action are the maps
\begin{equation}\label{eq:gtoug}
 G\ni g\mapsto u_g\in U
\end{equation}
satisfying the cocycle condition
\begin{equation}\label{eq:cocyccond}
 u_{gh} = u_g\cdot {}^g u_h,\qquad \forall g,h\in G.
\end{equation}
The space of all such cocycles is denoted by $Z^1(G,U)$. Two cocycles $(u_g)_g$ and $(v_g)_g$ are declared {\it equivalent} or {\it cohomologous} if there is some $w\in U$ with
\begin{equation*}
 v_g = w^{-1}\cdot u_g\cdot {}^g w,\qquad \forall g\in G,
\end{equation*}
and the set $H^1(G,U)$ of equivalence classes is the {\it first cohomology set} of $G$ with values in $U$.

These objects appear often in the theory of operator algebras (e.g., \cite[Definition X.1.4]{tak2}), in the process of deforming or ``twisting'' group actions on said algebras.

A smattering of other symbols featuring below:
\begin{itemize}\itemsep=0pt
\item $C^n(G,U):=C(G^n,U)$ is simply the space of maps $G^n\to U$ (which maps we refer to as {\it $n$-cochains}). It admits an obvious analytic-manifold structure, topology, etc., and is a~group when $U$ is abelian.

 Suppose $U$ is abelian. As explained in \cite[Section~VII.3]{ser-lf}, the space $C^n(G,U)$ can also be identified with that of $G$-module maps $G^{n+1}\to U$ (with $G$ acting on the domain diagonally, by left multiplication).
\item We have differentials
\[C^n(G,U)\xrightarrow{\quad\delta^n\quad}C^{n+1}(G,U),\] which compute the {\it cohomology} $H^n(G,U)$ as $
 H^n(G,U) = \ker \delta^n/\im \delta^{n-1}$.

\item $Z^n(G,U)\subset C^n(G,U)$ is the space of {\it $n$-cocycles} (with $n=0,1$ if $U$ is non-abelian). For abelian $U$, this is simply $Z^n(G,U):=\ker \delta^n$.
\item When $U$ is abelian,
 \begin{equation*}
 B^n(G,U):=\im\delta^{n-1}\subseteq Z^n(G,U)
 \end{equation*}
 is the space of {\it $n$-coboundaries}. We thus have
 \begin{equation*}
 H^n(G,U) = Z^n(G,U)/B^n(G,U).
 \end{equation*}
\end{itemize}
In all of the above, maps $G^n\to U$ are assumed continuous when working with compact $G$.

%%%%%%%%%%%%%%%%%%%%%%%%%%%%%%%%%%%%%%%%%%%%%%%%%%%%%%%%%%%%%%%%%%%%%%%%%%%%%
\subsection{Main results}\label{subse:gpmain}

\begin{Lemma}\label{le:cansplit}
 Let $G$ be a compact group acting linearly and continuously on a Banach space $E$ via
 \begin{equation*}
 G\times E\ni (g,e)\mapsto g\triangleright e\in E.
 \end{equation*}
 For every $n\in \bZ_{\ge 0}$, the map
 \begin{equation*}%\label{eq:ctobsurj}
 C^n(G,E)\xrightarrow{\quad\delta^n\quad} Z^{n+1}(G,E)
 \end{equation*}
 is a split surjection. Furthermore, we can choose splittings of norm $\le \sup_{g\in G}\|g\triangleright\|$.
\end{Lemma}
\begin{proof}
 The procedure is well known, and hinges on the ability to average against the Haar probability measure $\mu$ of $G$. In a slightly different context, the construction is described in \cite[Proof of Theorem~2.3]{moo-12} (which in turn cites earlier sources): given an $(n+1)$-cocycle ${a\in Z^{n+1}(G,E)}$,
 \begin{equation}\label{eq:splitgpcoc}
 b(g_1,\dots,g_{n}):=\int_{G}g^{-1}\triangleright a(g,\ g_1,\dots,g_n)\, \mathrm{d}\mu(g)
 \end{equation}
 is an $n$-cochain whose coboundary is precisely $a$.
\end{proof}

In the statement of Theorem \ref{th:cocyle-man-1} we will refer to the process of {\it twisting} an action by a~cocycle (e.g., \cite[Section~I.5.3, Example 2 preceding Proposition 34]{ser-gal} or \cite[Definition X.1.4, equation~(8)]{tak2}): if \eqref{eq:gactu} is an action and
\begin{equation*}
 u = (u_g)_g\in Z^1(G,U)
\end{equation*}
a 1-cocycle for it, one can define a new action by
\begin{equation*}
 g\triangleright x := u_g \cdot {}^g x \cdot u_g^{-1},\qquad g\in G,\quad x\in U.
\end{equation*}
Theorem \ref{th:cocyle-man-1} refers to the Lie-algebra version of this instead; see also Lemma \ref{le:equivact} for a more elaborate discussion in a slightly extended setup.

\begin{Theorem}\label{th:cocyle-man-1}
 Let $G$ be a compact group acting on a Banach Lie group $U$ with Lie algebra $\fu$.
 \begin{enumerate}[$(1)$]\itemsep=0pt
 \item\label{item:cocyle-man-1} The space
 \begin{equation*}
 Z^1(G,U)\subset C(G,U)
 \end{equation*}
 of $1$-cocycles is a closed analytic submanifold in the sense of {\rm \cite[\emph{Section}~5.8.3]{bourb-vars-17}}.

 \item\label{item:cocyle-man-1-tgz1} The tangent space to that subvariety at a cocycle $u\in Z^1(G,U)$ is $\rho_{u} Z^1(G,\fu)$, where
 \begin{itemize}\itemsep=0pt
 \item $\rho_{u}$ indicates right translation of tangent vectors by ${u}\in C(G,U)$;
 \item and the action `$\triangleright$' of $G$ on $\fu$ is the original one, twisted by the cocycle ${u}$:
 \begin{equation}\label{eq:defact}
 g\triangleright X:=\operatorname{Ad}_{u_g} {}^g X,\qquad \forall g\in G,\quad X\in \fu,
 \end{equation}
 with `$\operatorname{Ad}$' denoting the adjoint action of $U$ on its Lie algebra.
 \end{itemize}

 \item\label{item:cocyle-man-1-h1disc} The cohomology set $H^1(G,U)$, equipped with its quotient topology, is discrete.

\item\label{item:cocyle-man-1-coboundsplit} Every cocycle $u=(u_g)_g$ has a neighborhood $\cU$ in $Z^1(G,U)$ admitting an analytic map
 \begin{equation*}
 \cU\ni \left({u'}=(u'_g)_g\right)\xmapsto{\quad} v_{u'}\in U
 \end{equation*}
 such that
 \begin{align*}
 v_{u} =1\qquad\text{and}\qquad
 u'_g =v_{u'}^{-1}\cdot u_g\cdot {}^gv_{u'},\qquad \forall g\in G,\quad u'\in \cU.
 \end{align*}
 \end{enumerate}
\end{Theorem}
\begin{proof}
 We address the claims sequentially.

 \emph{Parts {\rm (\ref{item:cocyle-man-1})} and {\rm(\ref{item:cocyle-man-1-tgz1})}:} Fix, throughout the discussion, a 1-cocycle
 \begin{equation*}
 u=(u_g)_{g\in G}\in Z^1(G,U)
 \end{equation*}
 and the resulting twisted action \eqref{eq:defact} of $G$ on the Lie algebra $\fu$.

 A sufficiently small neighborhood $\cV$ of $u$ in the Lie group $C(G,U)$ consists of elements of the form
 \begin{equation*}
 v=(v_g)_g,\qquad v_g = \exp(X_g)\cdot u_g
 \end{equation*}
 for small-norm elements $X_g\in \fu$, with $\exp\colon\fu\to U$ the {\it exponential} (or {\it exponential function} \cite[discussion following Definition IV.1]{neeb-inf} ) of the Lie group.


 The additional condition that $v$ itself be a 1-cocycle reads
 \begin{equation*}
 \exp(X_{gh})u_{gh} = \exp(X_g)u_g \cdot \exp\left({}^gX_h\right){}^gu_h.
 \end{equation*}
 Or:
 \begin{align*}
 \exp(X_{gh})u_{gh} &= \exp(X_g)\cdot u_g\exp({}^gX_h)u_g^{-1}\cdot u_g\cdot{}^gu_h\\
 &= \exp(X_g)\cdot u_g\exp({}^gX_h)u_g^{-1}\cdot u_{gh}
 \quad\text{by \eqref{eq:cocyccond}}\\
 &=\exp(X_g)\cdot \exp(\operatorname{Ad}_{u_g}\,{}^gX_h)\cdot u_{gh}
 \quad\text{by the definition of the adjoint action}\\
 &=\exp(X_g)\cdot \exp(g\triangleright X_h)\cdot u_{gh}
 \quad\text{by the definition of the action `$\triangleright$'}.
 \end{align*}
 The rightmost factor cancels out to give
 \begin{equation*}
 \exp(X_{gh}) = \exp(X_g)\cdot \exp(g\triangleright X_h).
 \end{equation*}
 By the {\it BCH} formula ({\it Baker--Campbell--Hausdorff}: \cite[Definition IV.1.3 and Corollary IV.1.10]{neeb-lc}), this translates to
 \begin{equation*}
 X_{gh} = X_g + g\triangleright X_h + \frac 12\left[X_g,g\triangleright X_h\right] + \cdots.
 \end{equation*}
 When all norms $\|g\triangleright X_h\|$ are sufficiently small (which we may as well assume), the right-hand partial sum starting with the summand $\frac 12[-,-]$ onward is $O(\|X_g\|\cdot \|g\triangleright X_h\|)$ in ``big O notation''~\cite{bigo}: its norm is dominated by $C\|X_g\|\cdot \|g\triangleright X_h\|$ for some fixed constant $C>0$.

 If $\|(X_g)_g\| = \max_{g\in G} \|X_g\|$ is of order $\varepsilon$,
 \begin{equation*}
 \left(X_{gh} - X_g - g\triangleright X_h\right)_{g,h}\in B^2(G,\fu)
 \end{equation*}
 is a 2-coboundary and hence a 2-cocycle of size $O(\varepsilon^2)$. It follows from Lemma \ref{le:cansplit} that we can recover it as the coboundary of a {\it smaller-norm} cochain:
 \begin{equation*}
 X_{gh} - X_g - g\triangleright X_h
 =
 Y_{gh} - Y_g - g\triangleright Y_h
 \end{equation*}
 with $\|g\triangleright Y_{\bullet}\|=O\big(\varepsilon^2\big)$, i.e., $(X_g-Y_g)_g\in Z^1(G,\fu)$.
 Furthermore, the same Lemma \ref{le:cansplit} also ensures that $(X_g)_g\mapsto (Y_g)_g$ can be chosen to be analytic.

 The fact that $(Y_g)_g$ is in $O\big(\varepsilon^2\big)$ whenever $(X_g)_g$ is of order $\varepsilon$ also ensures that this map can be inverted, so we have an analytic chart on $\cU:=\cV\cap Z^1(G,U)$:
 \begin{equation*}
 \cU\ni (\exp(X_g)\cdot u_g)_g
 \xmapsto{\quad}
 (X_g)_g
 \xmapsto{\quad}
 (X_g-Y_g)\in Z^1(G,\fu),
 \end{equation*}
 onto some neighborhood of $0\in Z^1(G,\fu)$. This suffices to prove (\ref{item:cocyle-man-1}) and (\ref{item:cocyle-man-1-tgz1}): the notion of subvariety in the statement requires that the tangent space
 \begin{equation*}
 T_u Z^1(G,U)\le T_u C(G,U)
 \end{equation*}
 split, which is indeed the case, once more by Lemma \ref{le:cansplit}.

\emph{ Parts {\rm(\ref{item:cocyle-man-1-h1disc})} and {\rm(\ref{item:cocyle-man-1-coboundsplit})}:} Fix $u=(u_g)_g\in Z^1(G,U)$.
 The differential of the map
 \begin{equation}\label{eq:noncommdiff}
 U\ni v\mapsto \big(v^{-1}\cdot u_g\cdot {}^g v\big)_g\in Z^1(G,U)
 \end{equation}
 at the origin $1\in U$ is nothing but the cochain differential $C^0(G,\fu)\to C^1(G,\fu)$ (for the twisted action `$\triangleright$' discussed above), so it surjects onto the tangent space
 \begin{equation*}
 T_u Z^1(G,U) = Z^1(G,\fu) = B^1(G,\fu)
 \end{equation*}
 (by part (\ref{item:cocyle-man-1-tgz1})) with split kernel (see Lemma \ref{le:cansplit}).

 This makes \eqref{eq:noncommdiff} a {\it submersion} in the sense of \cite[Section~5.9.1\,(i)]{bourb-vars-17}, so that it also satisfies the other, equivalent conditions of \cite[Section~5.9.1]{bourb-vars-17}:
 \begin{itemize}\itemsep=0pt
 \item By \cite[Section~5.9.1\,(iv)]{bourb-vars-17} some neighborhood of $1\in U$ maps, through \eqref{eq:noncommdiff}, {\it onto} some open neighborhood of $u\in Z^1(G,U)$. This means that cocycles sufficiently close to $u$ are cohomologous to it, hence the discreteness of the cohomology set. This proves part (\ref{item:cocyle-man-1-h1disc}).
 \item And, also by \cite[Section~5.9.1\,(iv)]{bourb-vars-17}, \eqref{eq:noncommdiff} admits a {\it local section}: precisely what (\ref{item:cocyle-man-1-coboundsplit}) requires.
 \end{itemize}
 This concludes the proof as a whole.
\end{proof}
